Chapter~\ref{chap:results_analysis} evaluated AttackDRO++ through four connected
evidence panels. The aggregate whitebox results first showed that the proposed
method gives modest average and held-out gains relative to Multi-AT while
leaving Worst(8) and the separate AutoAttack checks broadly unchanged. The
per-attack and class-wise drilldowns then clarified where this behavior comes
from: the method acts less like a single-norm specialist and more like a
balanced multi-attack model, although hard animal-class cells under strong
\(\ell_\infty\) attacks remain persistent failure cases.

The graybox analysis tested whether the same pattern survives transferred
attacks from separately trained surrogate checkpoints. It supported the
whitebox interpretation while also showing that transfer accuracy must be read
carefully because specialist models can appear strong when attacks transfer
poorly across threat geometries. Finally, the ablation section examined whether
the observed behavior depends on key design choices, treating cluster count,
anchor strength, and refresh frequency as sensitivity checks rather than
standalone proof.

Overall, the supported conclusion is conservative: within the CIFAR-10 and
ResNet-18 setting studied here, AttackDRO++ improves cross-attack balance over
the uniform multi-attack baseline without establishing superiority on every
attack or extending robustness claims beyond the evaluated protocol. Chapter~\ref{chap:conclusion}
summarizes these findings, states the limitations, and outlines future
extensions.
